\documentclass[a4paper]{article}

\usepackage[fontset=fandol]{ctex}
\usepackage{amsmath, amssymb}
\usepackage{graphicx}
\usepackage[margin=2.5cm]{geometry}
\usepackage[most]{tcolorbox}
\usepackage{etoolbox}
\usepackage{listings}
\usepackage{hyperref}
\usepackage{booktabs}
\usepackage{subcaption}
\usepackage{float}
\usepackage{tikz}
\IfFileExists{pgfplots.sty}{
    \usepackage{pgfplots}
    \pgfplotsset{compat=1.18}
}{}

\newtcolorbox{knowledgebox}[1]{
    enhanced, colback=blue!5!white, colframe=blue!75!black, colbacktitle=blue!75!black,
    coltitle=white, fonttitle=\bfseries, title=#1, attach boxed title to top left={yshift=-2mm, xshift=2mm},
    boxrule=1pt, sharp corners
}

\newtcolorbox{importantbox}[1]{
    enhanced, colback=yellow!10!white, colframe=yellow!80!black, colbacktitle=yellow!80!black,
    coltitle=black, fonttitle=\bfseries, title=#1, sharp corners
}

\newtcolorbox{warningbox}[1]{
    enhanced, colback=red!5!white, colframe=red!75!black, colbacktitle=red!75!black,
    coltitle=white, fonttitle=\bfseries, title=#1, sharp corners
}

\lstset{
    language=Python,
    basicstyle=\ttfamily\small,
    keywordstyle=\color{blue},
    stringstyle=\color{red!60!black},
    commentstyle=\color{green!60!black},
    breaklines=true,
    frame=single,
    numbers=left,
    numberstyle=\tiny\color{gray},
    captionpos=b,
    extendedchars=false
}

\newcommand{\notetitle}{Appendix A: Spring 2026 最新课程主线与研究重排}
\newcommand{\noteauthors}{Codex}
\newcommand{\notedate}{2026-04-16}
\newcommand{\videochannel}{Robots That Learn / UC Berkeley}
\newcommand{\videopublishdate}{Spring 2026 course page snapshot}
\newcommand{\videoduration}{N/A}
\newcommand{\videourl}{https://robots-that-learn.github.io/}
\newcommand{\videocoverpath}{cover.jpg}

\begin{document}

\begin{titlepage}
\centering
{\Large Fall 2024 主书补充章\par}
\vspace{1.1cm}
{\huge\bfseries \notetitle\par}
\vspace{0.7cm}
{\large \noteauthors\par}
\vspace{0.3cm}
{\large \notedate\par}
\vspace{1.1cm}

\ifdefempty{\videocoverpath}{
{\small This appendix is reconstructed from the official Spring 2026 course page and reading list.\par}
}{
\includegraphics[width=0.82\textwidth,height=0.45\textheight,keepaspectratio]{\videocoverpath}\par
}

\vfill
\begin{tcolorbox}[width=0.9\textwidth, colback=black!2!white, colframe=black!60, sharp corners]
\textbf{来源边界}：本章明确补写 Fall 2024 主书的补充层；由于本次目录内没有公开视频，本章依据 Spring 2026 官方课程页、reading list、course notes excerpt、transcript.jsonl 与 slides.jsonl 重建。\par
\textbf{课程页面}：\href{\videourl}{\nolinkurl{\videourl}}
\end{tcolorbox}
\end{titlepage}

\tableofcontents
\newpage

\begin{figure}[H]
\centering
\includegraphics[width=0.78\textwidth]{overview_cover.jpg}
\caption{课程官方材料中的代表性页面，用于帮助读者在进入本讲正文前建立整体上下文。}
\end{figure}
\begin{figure}[H]
\centering
\includegraphics[width=0.78\textwidth]{cover.jpg}
\caption{Spring 2026 官方课程页在本补充章中的 provenance card。它并不替代正文，而是说明本章依据官方 course page 和 reading list 重构。}
\end{figure}

\section{本章定位与阅读方式}

本章是 Fall 2024 主书的补充章，不是新的主线讲义。它的职责不是重新讲一遍整门课，而是把 Spring 2026 官方课程页中已经发生的课程重排、阅读更新与研究重心迁移，整理成一个可直接并入主书的 appendix。

\subsection{为什么要单独做补充章}

Spring 2026 仍然以 robot learning 为主轴，但官方课表已经明显把课程组织方式改写了：biomechanics、kinematics、dexterous manipulation、reinforcement learning、behavior cloning、visual imitation、locomotion、navigation 和 language-based planning 这些主题仍在，但顺序、分组和 reading 选择都更偏向“最新研究前沿”而不是“Fall 2024 的历史编排”。

\begin{warningbox}{本章的证据边界}
本目录中没有公开视频可用，因此正文中的新增事实只来自 Spring 2026 官方课程页、reading list、course notes excerpt 以及本目录的结构化 transcript/slides 记录。凡是超出这些来源的内容，都只作为教材式延伸解释，不作为课程事实断言。
\end{warningbox}

\subsection{如何读这章}

如果你已经读过主书，可以把这一章当作两个作用不同的东西：
\begin{itemize}
  \item 第一，它是 Spring 2026 官方主线的“课程地图”；
  \item 第二，它是主书后续章节在写作时应吸收的“研究更新索引”。
\end{itemize}

换句话说，这一章不是为了替代主书，而是为了告诉读者主书在 2026 年应当怎么继续演化。

\section{Spring 2026 的课程主线重排}

Spring 2026 的 course page 仍然围绕 robot learning 展开，但它把课程更明显地分成了几块：从基础运动生物力学与机构学，到手与触觉，再到生成式动作建模、强化学习、模仿学习、腿足运动、导航，以及长程规划和语言。

\subsection{主线从“主题串联”变成“研究谱系”}

Fall 2024 更像是按主题顺序推进的一条学习路径；Spring 2026 则更像是一份研究谱系图。它保留了课程的教学目标，但把每个模块都放在更接近最近论文和工具链的位置上。

\begin{table}[H]
\centering
\small
\begin{tabular}{@{}p{0.28\textwidth}p{0.24\textwidth}p{0.38\textwidth}@{}}
\toprule
模块 & Spring 2026 讲次 & 课程意义 \\
\midrule
基础生物力学 / 机构学 & Lecture 2A / 2B & 先把人体与机器人运动学打底 \\
手、触觉与灵巧操作 & Lecture 4A / 4B & 把 manipulation 的感知接口讲完整 \\
发展、控制与神经视角 & Lecture 5A / 5B / 6A & 解释为何学习与控制必须并行 \\
生成式动作与 world models & Lecture 3 / 6B & 体现课程的最新研究增量 \\
RL / BC / visual imitation & Lecture 7AB / 8A / 8B & 维持主书中最关键的学习路线 \\
locomotion / navigation / language & Lecture 9AB / 10AB / 12AB & 把任务层级推向长程与语义 \\
\bottomrule
\end{tabular}
\caption{Spring 2026 官方课表相对 Fall 2024 的主题重排。}
\end{table}

\subsection{Lecture 1 到 Lecture 12 的组织方式}

从官方课程页看，Lecture 1 是 Introduction；Lecture 2A 是 Biomechanics of walking and running；Lecture 2B 是 Robot mechanisms - kinematics and dynamics；Lecture 3 是 Introduction to Diffusion Models；Lecture 4A/4B 继续围绕 hand、touch 与 proprioception；Lecture 5A/5B 进入 developmental perspective 与 robot dynamics/control/motion planning；Lecture 6A/6B 则把 computational neuroscience 和 Video World Models 放到一起。

后续的 Lecture 7AB 到 Lecture 10AB 依次覆盖 Reinforcement Learning、Behavior cloning、Visual Imitation、Case Studies in Locomotion、Case Studies in Navigation；Lecture 11AB 和 12AB 则把课程推进到 dexterous manipulation 和 long-horizon planning/language。

\subsection{本章小结}

Spring 2026 没有改变这门课的核心使命，但它改变了课程的“知识排序”。新的排序告诉读者：机器人学习不再只是“示范、控制、规划”的线性串联，而是一个从物理建模、生成式表示到长程语义规划的连续谱系。

\section{新增主题：Diffusion Models 与 Video World Models}

Spring 2026 最明显的新增就是 Lecture 3 的 Introduction to Diffusion Models，以及 Lecture 6B 的 Video World Models。它们并不是孤立插入的热门话题，而是课程主线向生成式建模与预测式表示学习扩展的标志。

\subsection{为什么 diffusion 会进入 robot learning}

扩散模型（diffusion models）最重要的不是“生成得像”，而是它提供了一种稳定建模多峰分布（multimodal distribution）的方式。对机器人来说，这一点尤其重要，因为一个观测往往对应多个合理动作。若直接用单峰回归去拟合，模型容易输出平均动作，反而更危险。

\begin{importantbox}{生成式策略的核心动机}
在机器人中，动作分布通常是多模态的。扩散模型、flow matching 和相关生成式方法的共同目标，都是避免把多个合理方案粗暴平均成一个不可靠动作。
\end{importantbox}

\subsection{Video World Models 的位置}

Lecture 6B 的 Video World Models 表明课程把“视频预测”与“控制”进一步打通。视频世界模型（video world model）可以理解为：系统不只看当前图像，还尝试预测未来视觉演化，从而为策略提供更长视野的中间表示。

在机器人学习中，这种表示有两层意义：
\begin{itemize}
  \item 它可以作为 planning 的中间层，帮助系统在潜在未来状态上做决策；
  \item 它也可以作为 imitation 的强化版，让动作学习不只对齐当前帧，而是对齐未来的可达后果。
\end{itemize}

\subsection{World Action Models 的信号}

官方 reading list 把 World Action Models are Zero-shot Policies 也纳入 Lecture 6B 的阅读，这说明课程已经把“world model”和“policy”之间的边界看得更紧了。世界动作模型（world action model）可以视为一种把环境预测和动作选择耦合起来的政策表示：它不是先预测再说，而是让预测本身服务于行动。

\subsection{Track2Act 的意义}

Track2Act: Predicting Point Tracks from Internet Videos enables Generalizable Robot Manipulation 指向另一条路线：从互联网视频中提取 point tracks，再把这些时空轨迹用于可泛化的操作学习。这里的关键不是“视频很大”，而是 point tracks 这种几何中间表示足够压缩，又保留了动作相关的空间结构。

\subsection{本章小结}

这一组新增主题说明，Spring 2026 已经不满足于“强化学习 + 模仿学习”的传统结构，而是把生成式建模、视频预测和动作决策看成同一条技术链上的不同环节。

\section{Reading 更新意味着什么}

Spring 2026 的 reading list 不只是“换几篇新论文”，而是在告诉读者：这门课的理论底座正在重写。Lecture 2B、3、6B、8A、8B、9AB、10AB 的 reading 变化尤其明显。

\subsection{几何和运动学的底座}

Lecture 2B 明确提醒学生先熟悉 3D rotations and translations 的表示，尤其是 exponential coordinates、twist 和 product of matrix exponentials。这里列出的 Li-Murray-Sastry textbook 和 Lynch-Park textbook，实际上把课程的 kinematics 基础重新拉回到了经典机器人学。

\begin{knowledgebox}{指数坐标（exponential coordinates）}
指数坐标（exponential coordinates）把旋转和刚体位姿统一到李群/李代数视角下：旋转矩阵可写成斜对称矩阵的指数，而带平移的刚体变换可写成 twist 的指数。对机器人前向运动学（forward kinematics）而言，这是一种更统一、更易组合的表示。
\end{knowledgebox}

\subsection{生成式建模的底座}

Lecture 3 的 Normalizing Flows for Probabilistic Modeling and Inference 与 Flow Matching for Generative Modeling 表明课程在生成式建模上选用的是两条非常清晰的理论线索。

\begin{itemize}
  \item Normalizing flows 强调可逆变换与显式似然；
  \item Flow matching 强调通过拟合速度场来构造连续生成过程。
\end{itemize}

这两条线都非常适合 robot learning，因为机器人动作并不总是单解问题。模型不仅要“生成一个动作”，还要能保留不确定性和多峰结构。

\subsection{操作、触觉与 locomotion 的阅读扩展}

Lecture 4A/4B 的 readings 把 human hand anatomy、A century of robotic hands、LEAP Hand、Human hand function 与 Touch 连接起来，说明这门课把 hand、tactile sensing 与 manipulation 视为同一问题域的不同层次。

Lecture 5A/5B 的 readings 则分别把 development of embodied cognition、internal models、control strategies in object manipulation tasks 和 visual locomotion 的跨模态监督拉进来。这一安排的含义是：机器人运动控制不只是 control theory 的问题，也不是纯学习的问题，而是 development、prediction 和 control 的交叉。

\subsection{学习、模仿和泛化}

Lecture 7AB、8A、8B 的 reading list 继续保留 Learning to Walk via Deep Reinforcement Learning、Learning Dexterous In-Hand Manipulation、Diffusion Policy 和 Universal Manipulation Interface。这说明课程并没有把 RL 与 imitation 视为过时路线，而是把它们当作大规模数据驱动控制的必要组成部分。

\subsection{本章小结}

Reading 更新的信号很直接：Spring 2026 并没有抛弃经典机器人学，而是在经典几何、控制和规划之上，引入更强的生成式建模、视频表征与跨模态学习。

\section{这对主书意味着什么}

如果把 Spring 2026 作为主书的补充层，那么它至少改变三件事。

\subsection{第一，课程不再是“学习算法单线叙事”}

主书不能只把 robot learning 写成 imitation、RL、BC 的线性发展史。Spring 2026 说明，几何建模、生成式建模、视频世界模型和语言规划已经共同构成 robot learning 的方法栈。

\subsection{第二，关节级动作和长程计划要被放在同一本书里讨论}

从手、触觉、locomotion 到 navigation 和 language，课程已经把“动作精度”与“任务跨度”放到同一个教学框架里。主书因此必须同时处理 low-level control 与 high-level planning。

\subsection{第三，补充章不应该只是附录注释}

Spring 2026 的内容已经足以作为主书后续版本的更新层。更合理的做法不是把它们压成一句“2026 有更新”，而是将它们独立成章，让读者能明确看见研究方向如何演化。

\section{本章小结}

本章把 Spring 2026 的官方课程页读成了一张“研究路线图”。它最重要的结论是：这门课依然是 robot learning，但它已经更明确地融合了生成式动作建模、视频世界模型、经典运动学、触觉操作和长程语义规划。

对主书来说，这意味着后续版本应该把 Spring 2026 视为“更新后的课程骨架”，而不是简单的补丁。

\section{总结与延伸}

这一章的价值不在于多列出了哪些 lecture，而在于它改变了读者对整门课的理解方式。Spring 2026 告诉我们，robot learning 正在从“会做动作”走向“会表示未来、会表达计划、会在生成式空间里做决策”。

如果把主书继续往前推进，最值得延伸的方向有三条：
\begin{itemize}
  \item 生成式策略与控制的统一：normalizing flows、diffusion、flow matching 如何成为 policy class；
  \item 视频世界模型与规划的统一：prediction 如何真正服务于 control；
  \item 语言、任务与几何的统一：language 如何和 manipulation、navigation、long-horizon planning 形成层级系统。
\end{itemize}

\subsection{本章小结}

Spring 2026 不是简单的课程重排，而是一次方法论更新。它提示主书在下一版中必须正面处理生成式建模、视频预测和长程规划，而不能只停留在“示范学习 + 强化学习”的旧框架里。

\end{document}
